\documentclass{article}
\usepackage[utf8]{inputenc}


\title{Trab ED2}
\author{Nomes}
\date{\today}




\begin{document}


\maketitle


\section{Introdução}
Este é um arquivo LaTeX criado no VS Code e versionado pelo GitHub!
Apresentar uma descrição geral do problema abordado, dos objetivos do projeto e dos algoritmos
implementados no sistema adaptativo, justificando brevemente as escolhas realizadas pelo grupo.
Acho que coloca o link do vídeo aqui


\section{Algoritmos Base}
    \subsection{Bubble Sort}
    explicação do Bubble


    \subsection{Insertion Sort}
    explicação do Insertion


    \subsection{Quick Sort}
    explicação do Quick


    \subsection{Heap Sort}
    explicação do Heap


    \subsection{Radix Sort}
    explicação do radix


\section{Arquitetura do Sistema}
Descrição:
• das métricas coletadas;
• das heurísticas;
• das decisões implementadas -> parte da bia


\section{Metodologia Experimental}
Descrever:
• ambiente utilizado;
• datasets;
• tamanhos testados;
• critérios de avaliação


\section{Resultados}
Apresentar tabelas, gráficos, métricas coletadas, com discussão crítica. Explicar:
• decisões corretas;
• decisões ruins;
• comportamentos inesperados;
• limitações observadas;
• impacto das características da entrada.
• impactos das entradas adversariais


\section{Reflexão sobre Uso de IA}
O grupo deverá informar:
• quais ferramentas de IA foram utilizadas;
• em quais etapas auxiliaram;
• erros encontrados nas respostas geradas;
• correções realizadas pelo grupo.


\section{Conclusão}
Resumo dos principais achados.


\section{Referências Bibliográficas}




\end{document}
